\section{Divisorial geometry of a remote entrance}\label{sec:v101-divisors}
We now identify the singularity datum behind a nonzero entrance order.
The datum is a pair: the ideal of the observation residual and the
ideal of the centre-path parameter, together with the accessible real
sectors. It is not the singularity of the remote image alone. The norm
on the residual supplies additional unit data for the leading coefficient.

Fix a remote representative and a sufficiently small compact isolating
neighbourhood $N$. Write
\[
 R(x,\delta)=f(x)-p(\delta),\qquad
 \mathfrak a=(R_1,\ldots,R_M),\qquad \mathfrak t=(\delta),\qquad
 d(\delta)=\min_{x\in N}\|R(x,\delta)\|.
\]
In this section $N$ is semialgebraic, $f$ is Nash on a neighbourhood
of $N$, and $p$ is Nash near zero. Assume that $R(x,0)=0$ has just
one solution in $N$, that it lies in the relative interior of $N$,
and that $d(\delta)>0$ for all sufficiently small positive $\delta$.
We also assume $p(0)=f(x_*)$ at that solution. This is the nonexact
alternative of Theorem~\ref{thm:singular-entrances} in the Nash category;
its broader continuous-semialgebraic statement is retained unchanged.
A Nash path after a positive ramification $\delta=s^q$ is treated by
applying the following construction to $s$ and dividing its order by $q$.

We use simultaneous real principalization of the \emph{finite} pair
$(\mathfrak a,\mathfrak t)$, with rectilinearization of the inequalities
of $N\times[0,\delta_0]$. The standard resolution and
rectilinearization results \cite{BM88,BM97} give, after shrinking and
covering a compact preimage, finitely many closed real sectors with
coordinates $(r,z)$, $r_i\ge0$, on which
\begin{equation}\label{eq:v101-monomial-pair}
 \delta=u(r,z)\prod_i r_i^{a_i},\qquad
 \|R\|=v(r,z)\prod_i r_i^{b_i},
 \qquad a_i,b_i\in\N,\qquad u,v>0.
\end{equation}
The units and their inverses are bounded on these sectors. To obtain
$v$, principalize the residual ideal, write its pulled-back generators
as $r^b U_1,\ldots,r^b U_M$ with no common zero, and take
$v=(\sum U_i^2)^{1/2}$. The charts parametrize the actual closed
feasible set, not the complex zero locus of its equations. Only sectors
in the closure of their positive-time interiors are retained.
We call their coordinate faces \emph{accessible real divisors}.
If a chart has a lower-dimensional source component, coordinates and
faces are taken in that component; identically zero-time components
are discarded. These operations are on the finite diagram just specified.

\begin{theorem}[Divisorial contact and the metric unit]\label{thm:v101-divisorial}
In any such finite real monomial presentation, an accessible index
with $a_i=0$ has $b_i=0$. Define
\begin{equation}\label{eq:v101-rho}
 \rho=\max_{\text{accessible }i,\ a_i>0}\frac{b_i}{a_i}.
\end{equation}
Then $\rho$ is positive and independent of the presentation. It has
the intrinsic real-arc characterization
\begin{equation}\label{eq:v101-arc-contact}
 \rho=\sup_{\gamma}
 \frac{\operatorname{ord}(\mathfrak a\circ\gamma)}
      {\operatorname{ord}(\delta\circ\gamma)},
\end{equation}
where $\gamma$ ranges over feasible positive-time semialgebraic Puiseux
arcs tending to $(x_*,0)$, and the order of an ideal is the minimum
order of its displayed generators. The supremum is attained by an arc
transverse to a suitable accessible divisor.

On each resolved sector the function $\delta^\rho/\|R\|$ extends
continuously to time zero as
\begin{equation}\label{eq:v101-unit-ratio}
 \Lambda(r,z)=\frac{u(r,z)^\rho}{v(r,z)}
                 \prod_i r_i^{\rho a_i-b_i}.
\end{equation}
Let $M$ be its maximum on the compact accessible zero-time preimage,
allowing all the sectors. Then $0<M<\infty$ and
\begin{equation}\label{eq:v101-unit-constant}
 d(\delta)=M^{-1}\delta^\rho+o(\delta^\rho).
\end{equation}
Equivalently, $\rho$ is the least real number $c$ for which
$\delta^c\le C\|R(x,\delta)\|$ throughout the feasible germ, for
some $C>0$. Thus the entrance exponent is a relative real
\L{}ojasiewicz contact exponent of $(\mathfrak a,\mathfrak t)$,
and the entrance coefficient is the reciprocal of a resolved unit maximum.
\end{theorem}

\begin{proof}
All assertions concern the finite real sectors of
\eqref{eq:v101-monomial-pair}. If $a_i=0<b_i$, set $r_i=0$ while
keeping the other boundary coordinates small and positive. The time
is positive and the residual vanishes. The image is feasible, which
contradicts $d(\delta)>0$. Thus $b_i=0$ whenever $a_i=0$. There are
positive-time points with both time and residual tending to zero,
for instance the fixed competitor $x_*$. Consequently at least one
accessible ratio $b_i/a_i$ is positive.

For $\rho$ in \eqref{eq:v101-rho}, every exponent
$\rho a_i-b_i$ is nonnegative. The bounded positive units therefore
make \eqref{eq:v101-unit-ratio} continuous and bounded on each compact
sector. Choose an index attaining the maximal ratio and approach a
generic point of its face, with all other boundary coordinates positive.
The ratio has a strictly positive value there. Its image at time zero
has zero residual and therefore is $(x_*,0)$. It follows that its
maximum $M$ over the zero-time preimage is positive and finite.

For an arbitrary feasible Puiseux arc, properness and semialgebraic
selection give a Puiseux lift after further ramification. On a sector
containing a final subarc, write $\beta_i=\operatorname{ord}r_i\ge0$;
coordinates tending to a nonzero value have order zero. Coordinates
with positive $a_i$ cannot vanish identically along a positive-time
arc. Coordinates with positive $b_i$ cannot vanish identically either,
since the residual is nonzero there. The units have order zero. Hence
\[
 \frac{\operatorname{ord}(\mathfrak a\circ\gamma)}
      {\operatorname{ord}(\delta\circ\gamma)}
 =\frac{\sum_i b_i\beta_i}{\sum_i a_i\beta_i}
 \le\max_{a_i>0}\frac{b_i}{a_i}.
\]
A transverse arc at a generic maximizing face gives equality. This
proves \eqref{eq:v101-arc-contact}. It proves independence of the
presentation without identifying a real sector with all real points
of a complex exceptional divisor.

For $\delta>0$, surjectivity onto the feasible set and compact attainment
give the exact identity
\begin{equation}\label{eq:v101-level-maximum}
 \max_{\text{resolved time }=\delta}\Lambda
                  =\frac{\delta^\rho}{d(\delta)}.
\end{equation}
Every sequence of maximizers on its left has a subsequence converging
to the compact zero-time preimage, so its upper limit is at most $M$.
For the reverse bound, take a zero-time point attaining $M$. It belongs
to the closure of the positive-time part. Curve selection supplies a
positive-time semialgebraic arc tending to it. Its time coordinate is
nonconstant and, after shortening the arc, strictly monotone, with
image containing every sufficiently small positive time. Evaluation
along this arc gives a lower limit at least $M$ in
\eqref{eq:v101-level-maximum}. This proves
\eqref{eq:v101-unit-constant}. Finally the bounded extension of
$\Lambda$ gives $\delta^\rho\le C\|R\|$. On a transverse maximizing
arc the ratio $\delta^c/\|R\|$ is unbounded for $c<\rho$.
This proves the last characterization.
\end{proof}

\begin{corollary}[Integral closure and the role of the norm]\label{cor:v101-integral}
Replacing the residual ideal by an ideal with the same integral closure
leaves \eqref{eq:v101-arc-contact} unchanged, with the same feasible
real sectors and time parameter. The leading coefficient need not be
preserved by this replacement. If the centre path extends $C^1$ to zero,
then $\rho\ge1$.
\end{corollary}
\begin{proof}
If $h$ is integral over an ideal $I$, an integral equation has the form
$h^n+a_1h^{n-1}+\cdots+a_n=0$, with $a_j\in I^j$. Along an arc,
$\operatorname{ord}h<\operatorname{ord}I$ would make $h^n$ the unique
lowest-order term, an impossibility. Ideals with the same integral
closure thus have the same minimum order on every arc, and
\eqref{eq:v101-arc-contact} applies. This is the valuative mechanism
behind the classical connection with \L{}ojasiewicz exponents
\cite{A2LJT08}. Multiplication of every residual generator by a positive
constant preserves the ideal but multiplies its norm and the entrance
coefficient by that constant. The latter requires the unit in
\eqref{eq:v101-unit-ratio}, not just the unnormed ideal.
Finally $x_*$ is an admissible competitor and
$\|R(x_*,\delta)\|=\|p(0)-p(\delta)\|=O(\delta)$.
Together with \eqref{eq:v101-unit-constant}, this forces $\rho\ge1$.
\end{proof}

\subsection{How the contact pair recovers the existing examples}
The singular realization in Theorem~\ref{thm:entrance-realization}
has, in probability coordinates, the residual
\[
 A(u^m-\delta)+Bu^n,\qquad n>m.
\]
The sector following the tangent contact can be parametrized by
\[
 \delta=u^m-u^n z=u^m(1-u^{n-m}z).
\]
There its two independent residual coordinates have common order $n$,
whereas time has order $m$. Thus $(a,b)=(m,n)$ and the accessible
contact ratio is $n/m$. No larger ratio is possible: if $u^m<\delta/2$
then $|u^m-\delta|\ge\delta/2$, whereas if $u^m\ge\delta/2$ then
$u^n\ge(\delta/2)^{n/m}$. Independence of $A,B$ and local equivalence
of the probability and Hellinger metrics give a lower bound by a
positive multiple of $\delta^{n/m}$ in both cases. In the tangent
sector the resolved Hellinger unit at $u=0$ is
\[
 \frac12\|zA+B\|_I.
\]
Its reciprocal maximum is $2/\|B\|_I$ when
$\langle A,B\rangle_I=0$, recovering the exact coefficient
$\|B\|_I/2$. This computation identifies the divisor and its metric
unit; it does not merely read a Puiseux exponent from an already
minimized scalar function.

In the regular regime of Proposition~\ref{prop:corner-criterion},
cone coercivity localizes minimizing displacements to order $\delta$.
The ordinary blow-up $x-x_*=\delta v$ has residual
$\delta(Jv-b)+O(\delta^2)$ on the feasible cone. Separation makes its
leading norm positive; the relevant contact is $(a,b)=(1,1)$, and the
unit optimization gives $\min_v\|Jv-b\|=\mu$. This recovers the first
coefficient of Theorem~\ref{thm:remote-general}. The signed second
coefficient and the exact equality-shell test still require the
second-order and exact feasible-chart arguments there. In particular,
\eqref{eq:v101-unit-constant} does not assign a side when a threshold
has the same leading term as $d(\delta)$.

For the two cubic identification walls, the direct feasible charts
prove precisely this regular localization, even where the channel
rank drops. Their entrance exponents are therefore one, while their
metric units and nonidentified-side openings differ. The divisorial
formula explains what the exponent measures; it does not replace the
exact cubic fibre classification, divide by a vanishing determinant,
or suppress either wall mechanism.
